\documentclass[11pt]{article}
\usepackage{setspace}
\usepackage{geometry}
\geometry{margin=1in}
\usepackage{natbib}
\bibliographystyle{apalike}

\setstretch{1.15}

\begin{document}

% --- Applicant Information ---
\noindent
\textbf{Zhongli Wang} \\
Phone: +49 160 9541 6015 \\
Email: nkwangzhongli@gmail.com \\
Website: https://sites.google.com/view/zhongli-wang \\
Bielefeld, Germany \\

\vspace{0.5cm}

% --- Recipient Information ---
\noindent
\textbf{TU Darmstadt} \\
Cluster of Excellence -- Reasonable Artificial Intelligence (RAI) \\
Darmstadt, Germany \\

\vspace{0.5cm}

% --- Subject Line ---
\noindent
\textbf{Application for Postdoctoral Position in Reasonable Artificial Intelligence} \\

\vspace{0.3cm}

\noindent Dear Members of the Hiring Committee,

I am writing to apply for the Postdoctoral Researcher position in the Cluster of Excellence \textit{Reasonable Artificial Intelligence (RAI)} at TU Darmstadt, specifically for the Systemic AI track. 
I am currently completing my Ph.D.\ in Economics at Bielefeld University and Università Cattolica del Sacro Cuore (EPOC Program, Marie Skłodowska-Curie Grant No.~956107). 
My research sits at the intersection of computational social science, economics, and AI systems, with a strong emphasis on building composable, robust, and accountable AI architectures for real-world deployment. 
I am particularly interested in how multi-agent systems can be designed with explicit reliability guarantees, modular composition principles, and transparent verification mechanisms—questions that align closely with RAI's mission to develop AI systems with human-like reasoning capabilities that generalize beyond narrow tasks.

My doctoral dissertation, \textit{Strategic Behavior and Economic Dynamics in ICT Markets: An Agent-based Framework}, develops computational models in which firms and consumers interact through networks, update beliefs under information frictions, and adapt pricing and investment decisions over time. 
Methodologically, the dissertation contributes three building blocks directly relevant to RAI: 
(i) multi-agent modeling of strategic interaction under uncertainty, 
(ii) scalable simulation and estimation using parallelized experiments, simulation-based inference, and neural-network surrogates, and 
(iii) reproducible research engineering with experiment tracking, robustness checks, and transparent documentation. 
I implement large-scale simulations in Python (with \texttt{ray} for parallelism) and use simulation-based econometrics to infer behavioral parameters and evaluate counterfactual policies—skills that translate directly to building and validating complex agentic AI systems.

In parallel, I have developed a second research stream on AI-enabled scientific workflows. 
In the working paper \textit{Agentic Workflows for Economic Research: Design and Implementation}, we propose a modular architecture of specialized LLM agents that collaborate across stages of the research lifecycle (data collection, coding, modeling, evaluation), with human-in-the-loop control, auditability, and versioning. 
The architecture builds on recent advances in tool-augmented language models, retrieval-augmented generation, and reasoning-action integration. 
This experience directly translates to RAI's emphasis on compositional systems, pipeline governance, and robust deployment: building AI systems from interoperable modules, enforcing reproducibility, and designing evaluation protocols that detect failure modes beyond benchmark performance.

My proposed research project, \textit{Composable Agentic Systems for Reliable Reasoning Pipelines in Real-World Text-Rich Environments}, addresses a critical challenge in modern AI deployment: how to build end-to-end pipelines that chain multiple components (retrieval, extraction, reasoning, verification, summarization) while maintaining reliability and robustness under domain shift, ambiguous inputs, and compounding errors. 
I will develop and evaluate a modular agentic architecture in which specialized AI modules (LLM-based and classical ML/NLP components) are composed via explicit interfaces, shared state representations, and verification routines. 
The project will focus on neurosymbolic hybrid composition, combining LLM modules with structured components (entity linking, rule-based constraints, symbolic checks), and will establish reliability metrics including calibration, selective prediction, trace completeness, and robustness under perturbations.

Beyond research, I am strongly motivated by the position's emphasis on helping coordinate and build one research lab. 
My doctoral work, supported by the \textit{Marie Skłodowska-Curie Actions} and the \textit{Bielefeld Young Researchers' Fund}, required designing and executing complex empirical and computational workflows across multiple institutions. 
I have experience developing research infrastructure, organizing workshops, maintaining reproducible workflows across collaborators, and managing collaborative projects. 
I would contribute actively to lab coordination, project acquisition, mentoring of junior researchers, and building the technical infrastructure needed to support RAI's ambitious research agenda.

RAI offers a uniquely appropriate environment to develop this research agenda. 
The cluster's focus on systemic foundations, multi-modality, adaptation, and cognitive challenges matches my interests in agentic architectures, robustness, and evaluation. 
I am particularly excited about contributing to RAI's mission of developing AI systems that are composable into larger software systems, robust to distribution shift and adversarial environments, and transparent enough to support verification, debugging, and responsible use. 
My background in multi-agent modeling, computational economics, and AI-enabled workflows positions me well to bridge theoretical foundations with practical system-building, while my experience with research coordination and infrastructure development will enable me to contribute meaningfully to lab management and collaborative research initiatives.

Thank you for considering my application. I would be honored to join TU Darmstadt's RAI cluster as a postdoctoral researcher and to contribute to the development of Reasonable AI as both an engineering and scientific discipline. 
I look forward to the possibility of working together.

\noindent Sincerely,\\[12mm]
\noindent Zhongli Wang

% \vspace{5mm}

% \bibliography{NFR_Template_Statements/references1}

\end{document}
